% Luc Destombe --- licence CC BY-NC-SA 4.0
% https://creativecommons.org/licenses/by-nc-sa/4.0/deed.fr
% Fichier autonome : compiler avec pdflatex, deux fois (signets, nombre de pages).
\def\lyceeversion{1.1}
\documentclass[11pt,a4paper]{article}
\makeatletter
\newif\ifeleve
\elevefalse

\newif\iflycee@cm
\newif\iflycee@palatino

\RequirePackage[utf8]{inputenc}
\RequirePackage[T1]{fontenc}
\RequirePackage[french]{babel}
\RequirePackage{etoolbox}

\newcommand{\ShorthandsOff}{\@ifundefined{shorthandoff}{}{\shorthandoff{;:!?}}}
\newcommand{\ShorthandsOn}{\@ifundefined{shorthandon}{}{\shorthandon{;:!?}}}
\ShorthandsOff
\AtBeginDocument{\ShorthandsOn}
\AtBeginEnvironment{tikzpicture}{\ShorthandsOff}

\RequirePackage{geometry}
\RequirePackage{fancyhdr}
\RequirePackage{lastpage}
\RequirePackage{multicol}
\RequirePackage{array}
\RequirePackage{enumitem}
\RequirePackage{xcolor}
\RequirePackage{needspace}

\RequirePackage{amsmath,amssymb}
\iflycee@cm\else\RequirePackage{fourier}\fi
\DeclareMathAlphabet{\mathcal}{OMS}{cmsy}{m}{n}
\renewcommand{\ttdefault}{lmtt}
\begingroup\catcode`\_=\active \catcode`\^=\active
\global\let_\sb \global\let^\sp
\endgroup
\AtBeginDocument{\catcode`\_=12 \mathcode`\_="8000
  \catcode`\^=12 \mathcode`\^="8000 }
\newlength{\retraitcalculs}
\setlength{\retraitcalculs}{2em}

\RequirePackage{tikz}
\usetikzlibrary{arrows.meta,calc,babel,shapes.misc}
\RequirePackage[most]{tcolorbox}
\tcbuselibrary{skins,breakable}

\definecolor{coulDef}{HTML}{1F4E79}
\definecolor{coulProp}{HTML}{7030A0}
\definecolor{coulRem}{HTML}{C55A11}
\definecolor{coulEx}{HTML}{2E7D32}
\definecolor{coulExo}{HTML}{B03A2E}
\definecolor{coulPart}{HTML}{1F4E79}
\definecolor{coulMeth}{HTML}{1B6E4F}
\definecolor{coulCourbe}{HTML}{0B5ED7}
\definecolor{coulCourbeB}{HTML}{C00000}
\definecolor{coulCourbeC}{HTML}{2E7D32}
\definecolor{coulGrille}{HTML}{8C8C8C}
\definecolor{coulRep}{HTML}{1B6E4F}
\definecolor{coulBar}{HTML}{2E7D32}

\newcommand{\R}{\mathbb{R}}
\newcommand{\Rs}{\mathbb{R}^{*}}
\newcommand{\N}{\mathbb{N}}
\newcommand{\Ns}{\mathbb{N}^{*}}
\newcommand{\Z}{\mathbb{Z}}
\newcommand{\Q}{\mathbb{Q}}
\renewcommand{\le}{\leqslant}
\renewcommand{\ge}{\geqslant}
\renewcommand{\approx}{\simeq}
\newcommand{\vect}[1]{\overrightarrow{#1}}
\newcommand{\norme}[1]{\left\lVert #1 \right\rVert}
\newcommand{\intervalle}[2]{\left[#1\,;#2\right]}
\RequirePackage{cancel}
\renewcommand{\CancelColor}{\color{coulExo}}
\newcommand{\fois}[1]{\times{\color{coulExo}#1}}
\newcommand{\barre}[1]{\cancel{#1}}
\newlength{\lycee@eqh}\newlength{\lycee@eqd}
\newsavebox{\lycee@eqa}\newsavebox{\lycee@eqb}\newsavebox{\lycee@eqc}
\newcommand{\lycee@eqcol}[1]{\hbox{$\renewcommand{\arraystretch}{1.35}%
  \begin{array}[t]{@{}>{\displaystyle}r@{}>{\displaystyle{}}l@{}}#1\end{array}$}}
\newcommand{\lycee@eqmesure}[1]{%
  \ifdim\ht#1>\lycee@eqh\setlength{\lycee@eqh}{\ht#1}\fi
  \ifdim\dp#1>\lycee@eqd\setlength{\lycee@eqd}{\dp#1}\fi}
\newcommand{\lycee@eqtrait}{\hspace{1.2em}ou\hspace{1.2em}}
\newcommand{\lycee@eqiff}{$\iff$\hspace{0.8em}}
\newcommand{\eqproduit}[2]{\par\smallskip\noindent C'est une équation produit :
  $A\times B=0 \iff A=0$ ou $B=0$.\par\nopagebreak\smallskip
  \begingroup\setlength{\lycee@eqh}{0pt}\setlength{\lycee@eqd}{0pt}%
  \setbox\lycee@eqa=\lycee@eqcol{#1}\setbox\lycee@eqb=\lycee@eqcol{#2}%
  \lycee@eqmesure\lycee@eqa\lycee@eqmesure\lycee@eqb%
  \noindent\hspace*{\retraitcalculs}\lycee@eqiff\usebox\lycee@eqa\lycee@eqtrait\usebox\lycee@eqb\endgroup\par\smallskip}
\newcommand{\eqproduittrois}[3]{\par\smallskip\noindent C'est une équation produit :
  $A\times B\times C=0 \iff A=0$ ou $B=0$ ou $C=0$.\par\nopagebreak\smallskip
  \begingroup\setlength{\lycee@eqh}{0pt}\setlength{\lycee@eqd}{0pt}%
  \setbox\lycee@eqa=\lycee@eqcol{#1}\setbox\lycee@eqb=\lycee@eqcol{#2}\setbox\lycee@eqc=\lycee@eqcol{#3}%
  \lycee@eqmesure\lycee@eqa\lycee@eqmesure\lycee@eqb\lycee@eqmesure\lycee@eqc%
  \noindent\hspace*{\retraitcalculs}\lycee@eqiff\usebox\lycee@eqa\lycee@eqtrait\usebox\lycee@eqb\lycee@eqtrait\usebox\lycee@eqc\endgroup
  \par\smallskip}

\tikzset{grille/.style={coulGrille,line width=0.4pt}}
\newcommand{\quadrillage}[4]{\draw[grille,step=1] (#1,#2) grid (#3,#4);}
\tikzset{
  croix/.style={cross out,draw,line width=0.6pt,inner sep=0pt,minimum size=4pt},
  croixdroite/.style={croix,rotate=45,line width=0.8pt},
}
\newcommand{\pt}[3]{\path (#1) node[croix]{} node[#2,font=\small,inner sep=2.5pt]{$#3$};}

\newcommand{\lycee@repere}[4]{%
  \draw[grille,step=1] (#1,#2) grid (#3,#4);
  \draw[-{Stealth[length=2mm]},thick] (#1,0)--({#3+0.4},0) node[below left=-1pt]{$x$};
  \draw[-{Stealth[length=2mm]},thick] (0,#2)--(0,{#4+0.4}) node[below left=-1pt]{$y$};
  \node[below left,font=\scriptsize,inner sep=1.5pt] at (0,0) {$0$};}
\newcommand{\lycee@gradx}[1]{\foreach \k/\l in {#1}{%
  \draw (\k,0.07)--(\k,-0.07) node[below,font=\scriptsize,inner sep=1.5pt]{$\l$};}}
\newcommand{\lycee@grady}[1]{\foreach \k/\l in {#1}{%
  \draw (0.07,\k)--(-0.07,\k) node[left,font=\scriptsize,inner sep=1.5pt]{$\l$};}}
\newcommand{\lycee@intff}[2]{\left[#1\,;#2\right]}
\newcommand{\lycee@intfo}[2]{\left[#1\,;#2\right[}
\newcommand{\lycee@intof}[2]{\left]#1\,;#2\right]}
\newcommand{\lycee@intoo}[2]{\left]#1\,;#2\right[}
\newcommand{\lycee@ens}[1]{\left\{#1\right\}}
\AtBeginDocument{%
  \providecommand{\repere}{\lycee@repere}%
  \providecommand{\gradx}{\lycee@gradx}%
  \providecommand{\grady}{\lycee@grady}%
  \providecommand{\Cf}{\mathcal{C}\sb{f}}%
  \providecommand{\Cg}{\mathcal{C}\sb{g}}%
  \providecommand{\intff}{\lycee@intff}%
  \providecommand{\intfo}{\lycee@intfo}%
  \providecommand{\intof}{\lycee@intof}%
  \providecommand{\intoo}{\lycee@intoo}%
  \providecommand{\ens}{\lycee@ens}}

\newlist{questions}{enumerate}{3}
\setlist[questions,1]{label=\arabic*\textdegree),leftmargin=*,itemsep=2pt,topsep=3pt}
\setlist[questions,2]{label=\alph*),leftmargin=*,itemsep=1pt,topsep=2pt}
\setlist[questions,3]{label=\roman*),leftmargin=*,itemsep=1pt,topsep=1pt}
\setlist[itemize]{label=\textbullet,leftmargin=*,itemsep=2pt,topsep=3pt}

\newcommand{\besoinplace}[1]{\par\penalty\@M\begingroup
  \dimen@=#1\relax\dimen@ii=\pagegoal\advance\dimen@ii-\pagetotal
  \ifdim\dimen@>\dimen@ii\ifdim\dimen@ii>\z@\vfil\fi\break\fi\endgroup}

\newcommand{\lycee@pied}{\thepage/\pageref{LastPage}}
\newcommand{\entete}[2]{%
  \pagestyle{fancy}\fancyhf{}%
  \fancyhead[L]{\footnotesize\color{coulPart}\textbf{#1}}%
  \fancyhead[R]{\footnotesize\color{coulPart}\textbf{#2}}%
  \fancyfoot[C]{\footnotesize\lycee@pied}%
  \renewcommand{\headrulewidth}{0.4pt}%
  \renewcommand{\headrule}{\hbox to\headwidth{%
    \color{coulPart!40}\leaders\hrule height \headrulewidth\hfill}}}

\RequirePackage{ccicons}
\newcommand{\lycee@licence}{{\scriptsize\color{gray}%
  \copyright\ Luc Destombe --- \ccbyncsa\ CC BY-NC-SA 4.0}}
\newif\iflycee@licence \lycee@licencetrue
\newcommand{\sanslicence}{\lycee@licencefalse}
\AtBeginDocument{\iflycee@licence\AddToHookNext{shipout/foreground}{%
  \put(0,0){\rlap{\kern\dimexpr1in+\hoffset+\oddsidemargin+\textwidth\relax
    \llap{\raisebox{-\dimexpr1in+\voffset+\topmargin+\headheight+\headsep
      +\textheight+\footskip\relax}{\lycee@licence}}}}}\fi}

\newcommand{\lycee@site}{faire-des-maths.fr}
\newcommand{\lycee@version}{{\scriptsize\color{gray}\lycee@site
  \ifdefined\lyceeversion\ --- Version~\lyceeversion\fi}}
\AtBeginDocument{\iflycee@licence\AddToHookNext{shipout/foreground}{%
  \put(0,0){\rlap{\kern\dimexpr1in+\hoffset+\oddsidemargin\relax
    \raisebox{-\dimexpr1in+\voffset+\topmargin+\headheight+\headsep
      +\textheight+\footskip\relax}{\lycee@version}}}}\fi}

\newcommand{\formulesserrees}{%
  \setlength{\abovedisplayskip}{3pt}\setlength{\belowdisplayskip}{3pt}%
  \setlength{\abovedisplayshortskip}{2pt}\setlength{\belowdisplayshortskip}{3pt}}

\setlength{\parindent}{0pt}

\RequirePackage[implicit=false,hidelinks,pdfencoding=auto,psdextra,bookmarksopen,
  bookmarksopenlevel=1,pdfcreator={},pdfproducer={}]{hyperref}
\RequirePackage{bookmark}
\hypersetup{pdfauthor={Luc Destombe},pdfkeywords={CC BY-NC-SA 4.0}}
\newcounter{lycee@signet}
\newcommand{\signet}[3][]{\stepcounter{lycee@signet}%
  \edef\lycee@dest{\if\relax\detokenize{#1}\relax
    signet.\arabic{lycee@signet}\else#1\fi}%
  \hypertarget{\lycee@dest}{}\bookmark[level=#2,dest=\lycee@dest]{#3}}
\pdfstringdefDisableCommands{%
  \def\R{R}\def\N{N}\def\Z{Z}\def\Q{Q}\def\vect#1{#1}\def\Cf{Cf}\def\Cg{Cg}\def\fois#1{×#1}%
  \def\barre#1{#1}\def\intervalle#1#2{[#1;#2]}\def\intff#1#2{[#1;#2]}%
  \def\textsuperscript#1{#1}\def\dfrac#1#2{#1/#2}\def\frac#1#2{#1/#2}%
  \def\vec#1{#1⃗}}%
\newcommand{\res}[1]{#1}
\geometry{top=2cm,bottom=1cm,left=1cm,right=1cm,headheight=14pt,footskip=0.6cm}

\RequirePackage{pgfplots}
\pgfplotsset{compat=1.18}
\RequirePackage{tkz-tab}

\newcounter{partiecours}
\renewcommand{\thepartiecours}{\Roman{partiecours}}
\newcommand{\partiecours}[1]{%
  \refstepcounter{partiecours}%
  \Needspace*{8\baselineskip}%
  \bigskip
  \noindent\begin{tcolorbox}[enhanced,colback=coulPart,colframe=coulPart,
      boxrule=0pt,arc=2pt,left=8pt,right=8pt,top=4pt,bottom=4pt,
      before skip=6pt,after skip=10pt]
    \color{white}\bfseries\large \thepartiecours{} --- #1%
    \signet[partie-\arabic{partiecours}]{1}{\thepartiecours{} --- #1}
  \end{tcolorbox}%
  \addcontentsline{toc}{section}{\thepartiecours{} --- #1}%
}

\newtcolorbox{boitecours}[3][]{%
  enhanced,breakable,
  colback=white,colframe=#2,coltitle=white,
  fonttitle=\bfseries,
  attach boxed title to top left={xshift=8pt,yshift=-\tcboxedtitleheight/2},
  boxed title style={colback=#2,arc=2pt,boxrule=0pt},
  boxrule=0.9pt,arc=2pt,
  left=8pt,right=8pt,top=8pt,bottom=6pt,
  title={#3},#1}

\newcounter{methode}
\newenvironment{definitionbox}[1][Définition]%
  {\begin{boitecours}[phantom={\signet{2}{#1}}]{coulDef}{#1}}{\end{boitecours}}
\newenvironment{proprietebox}[1][Propriété]%
  {\begin{boitecours}[phantom={\signet{2}{#1}}]{coulProp}{#1}}{\end{boitecours}}
\newenvironment{methodebox}[1][Méthode]{\stepcounter{methode}%
  \begin{boitecours}[phantom={\signet[methode-\arabic{methode}]{2}{#1}}]{coulMeth}{#1}}%
  {\end{boitecours}}

\newcommand{\etape}[1]{\par\smallskip\textbf{\color{coulMeth}#1}\par\nobreak\smallskip}
\newcommand{\enonce}[1]{\emph{#1}\par\smallskip}
\newcommand{\reponse}[1]{\par\smallskip
  {\footnotesize\itshape\color{black!55}Réponses : #1}}
\newcommand{\demo}[1]{\textbf{\color{coulMeth}Démonstration.} #1}
\newcommand{\afaire}[1]{%
  \par\smallskip
  \noindent{\small\color{coulExo}$\blacktriangleright$\ \textbf{À faire :}
    fiche d'exercices du chapitre, #1.}\par\smallskip}

\newenvironment{calculs}{\setlength{\jot}{5pt}\@fleqntrue\@mathmargin\retraitcalculs
  \setlength{\abovedisplayskip}{4pt}\setlength{\belowdisplayskip}{4pt}%
  \setlength{\abovedisplayshortskip}{4pt}\setlength{\belowdisplayshortskip}{4pt}%
  \csname alignat*\endcsname{2}}%
  {\csname endalignat*\endcsname}
\newcommand{\rem}[1]{\text{\small\itshape\color{black!60}#1}}
\newcommand{\explique}[2]{\par\smallskip\noindent
  \begin{minipage}[t]{0.57\linewidth}\raggedright #1\end{minipage}\hfill
  \begin{minipage}[t]{0.40\linewidth}\raggedright\leavevmode
    {\small\itshape\color{black!60}#2}\end{minipage}\par}
\newcommand{\cas}[1]{\par\smallskip\noindent\textbf{\color{coulExo}#1}\vadjust{\nobreak}}
\newcommand{\calc}[1]{\par\smallskip\textbf{\color{coulExo}#1}\par\nobreak}

\newtcolorbox{remarquebox}[1][Remarque]{%
  enhanced,breakable,colback=white,colframe=coulRem,
  boxrule=0pt,leftrule=2.5pt,arc=0pt,outer arc=0pt,
  left=8pt,right=6pt,top=5pt,bottom=5pt,
  before upper={\textbf{\color{coulRem}#1}\par\smallskip}}

\newtcolorbox{exemplebox}[1][Exemple]{%
  enhanced,breakable,colback=white,colframe=coulEx,
  boxrule=0pt,leftrule=2.5pt,arc=0pt,outer arc=0pt,
  left=8pt,right=6pt,top=5pt,bottom=5pt,
  before upper={\textbf{\color{coulEx}#1}\par\smallskip}}

\pgfplotsset{
  repere/.style={
    axis lines=middle,
    axis line style={-{Stealth[length=2.4mm]},thick},
    every axis x label/.style={at={(ticklabel* cs:1.0)},anchor=west},
    every axis y label/.style={at={(ticklabel* cs:1.0)},anchor=south},
    label style={font=\footnotesize},
    tick label style={font=\scriptsize},
    grid=both,
    major grid style={grille},
    minor tick num=0,
    tick align=inside,
    clip mode=individual,
  },
}
\tikzset{
  courbe/.style={coulCourbe,line width=1.1pt,smooth},
  courbeB/.style={coulCourbeB,line width=1.1pt,smooth},
  courbeC/.style={coulCourbeC,line width=1.1pt,smooth},
}

\newlist{etapes}{enumerate}{1}
\setlist[etapes]{label=\textbf{\arabic*.},leftmargin=*,itemsep=1pt,topsep=2pt}

\newlist{expressions}{enumerate}{1}
\setlist[expressions]{label={\Alph*\ $=$},leftmargin=*,itemsep=3pt,topsep=3pt}

\newcounter{exo}
\renewcommand{\theexo}{\ifnum\value{exo}<10 0\fi\arabic{exo}}
\newtcolorbox{exercice}[1][]{%
  enhanced,breakable,
  colback=white,colframe=coulExo,
  boxrule=0.9pt,arc=2pt,
  left=8pt,right=8pt,top=8pt,bottom=6pt,
  coltitle=white,fonttitle=\bfseries,
  attach boxed title to top left={xshift=8pt,yshift=-\tcboxedtitleheight/2},
  boxed title style={colback=coulExo,arc=2pt,boxrule=0pt},
  phantom={\signet{2}{Exercice \arabic{exo}}},
  title={Exercice \theexo\ifx\relax#1\relax\else{} --- #1\fi}}
\newenvironment{exo}[1][\relax]{\refstepcounter{exo}\begin{exercice}[#1]}{\end{exercice}}

  \newenvironment{correction}{\par}{\par}
  \newcommand{\autableau}{}
  \newcommand{\etapeexemples}{\etape{Exemples corrigés}}
  \newcommand{\etapetableau}[1]{\etape{#1}}
  \newenvironment{exempletableau}[1][Exemple]%
    {\begin{exemplebox}[#1]}{\end{exemplebox}}

\RequirePackage{comment}
\excludecomment{correctionavj}

\newcommand{\coursEntete}{}
\newcommand{\coursNiveau}{}
\newcommand{\coursTitre}{}
\newcommand{\coursSurTitre}{}
\newcommand{\coursSousTitre}{}

\newcommand{\enteteCours}[1]{\renewcommand{\coursEntete}{#1}}
\newcommand{\chapitrecours}[2]{\enteteCours{Chapitre #1 --- #2}}
\newcommand{\niveaucours}[1]{\renewcommand{\coursNiveau}{#1}}
\newcommand{\titrecours}[3][]{%
  \renewcommand{\coursTitre}{#2}%
  \renewcommand{\coursSurTitre}{#1}%
  \renewcommand{\coursSousTitre}{#3}}

\pagestyle{fancy}
\fancyhf{}
\fancyhead[L]{\footnotesize\color{coulPart}\textbf{\coursEntete}}
\fancyhead[R]{\footnotesize\color{coulPart}\textbf{\coursNiveau}}
\fancyfoot[C]{\footnotesize\thepage}
\renewcommand{\headrulewidth}{0.4pt}
\renewcommand{\headrule}{\hbox to\headwidth{\color{coulPart!40}\leaders\hrule height \headrulewidth\hfill}}

\setlength{\parskip}{2pt}

\newcommand{\lycee@mentioneleve}{\\[2pt]{\small\itshape Version élève}}
\AtBeginDocument{%
  \begin{center}
    \begin{tcolorbox}[enhanced,width=0.72\linewidth,colback=white,colframe=coulPart,
        boxrule=1.6pt,arc=3pt,halign=center,top=10pt,bottom=10pt]
      \ifx\coursSurTitre\empty
        {\Huge\bfseries\color{coulPart} \coursTitre}\\[4pt]
      \else
        {\Huge\bfseries\color{coulPart} \coursTitre}\\[3pt]
        {\Large\bfseries\color{coulPart} \coursSurTitre}\\[5pt]
      \fi
      {\normalsize \coursSousTitre}%
      
    \end{tcolorbox}
  \end{center}%
}

\RequirePackage{tikz}
\RequirePackage{xcolor}

\usetikzlibrary{decorations.text}
\definecolor{lycee@metal}{HTML}{1F4E79}
\definecolor{lycee@manche}{HTML}{8A5A2B}
\definecolor{lycee@etincelle}{HTML}{F28C28}
\newcommand{\lycee@pluriel}[2]{#1 #2\ifnum#1>1 s\fi}
\newcommand{\lycee@sceau}{%
  \edef\lycee@bas{\lycee@pluriel{\lyceerelectures}{relecture} et
    \lycee@pluriel{\lyceerodages}{rodage}}%
  \begin{tikzpicture}[line join=round]
    \fill[coulPart!10,draw=coulPart,line width=.8pt] (0:2.2)
      \foreach \i in {1,...,71}{ -- (\i*5:{isodd(\i)?2.1:2.2}) } -- cycle;
    \draw[coulPart,line width=1.2pt] (0,0) circle (1.93);
    \draw[coulPart,line width=.4pt] (0,0) circle (1.85);
    \draw[coulPart,line width=.4pt] (0,0) circle (1.2);
    \path[decorate,decoration={text along path,text align=center,
        text={|\fontsize{9.5}{11}\selectfont\bfseries\color{coulPart}|TEST{\'E} AVEC MES {\'E}L{\`E}VES}}]
      (180:1.45) arc (180:0:1.45);
    \path[decorate,decoration={text along path,text align=center,
        text/.expanded={|\noexpand\fontsize{8.5}{10}\noexpand\selectfont\noexpand\color{coulPart}|\lycee@bas}}]
      (195:1.68) arc (195:345:1.68);
    \begin{scope}[x=.25mm,y=-.25mm]
      \begin{scope}[shift={(-32,-40)}]%
        \fill[lycee@metal] (23,31) rectangle (41,49);
        \draw[lycee@etincelle,line width=.5pt,line cap=round]
          (20,30)--(15,27) (19,34)--(14,35) (44,30)--(49,27) (45,34)--(50,35);
        \fill[lycee@manche,rounded corners=.3mm] (38,24.6) rectangle (60,27.6);
        \fill[lycee@metal,rounded corners=.4mm] (25,20) rectangle (39,31);
        \fill[lycee@etincelle,rounded corners=.2mm] (25,20) rectangle (28,31);
      \end{scope}
    \end{scope}
    \ifdefined\lyceesceauplusieurs\def\lycee@annee{Ann{\'e}es~:}\else
      \def\lycee@annee{Ann{\'e}e}\fi
    \path[decorate,decoration={text along path,text align=center,
        text/.expanded={|\noexpand\fontsize{8}{10}\noexpand\selectfont\noexpand\color{coulPart}|%
          \unexpanded\expandafter{\lycee@annee} \unexpanded\expandafter{\lyceesceauannees}}}]
      (95:1.08) arc (95:445:1.08);
  \end{tikzpicture}}
\AtEndDocument{\ifdefined\lyceesceauannees
  \par\besoinplace{5cm}\vspace{1.5em}%
  \begin{center}\lycee@sceau\end{center}\fi}
\makeatother

\chapitrecours{2}{Limites de suites}
\niveaucours{T\textsuperscript{le} spécialité}
\titrecours{LIMITES DE SUITES}{Cours --- Terminale spécialité}

\begin{document}

\newtcblisting{python}{%
  listing only,enhanced,breakable,
  colback=black!3,colframe=black!35,boxrule=0.5pt,arc=1pt,
  left=6pt,right=6pt,top=2pt,bottom=2pt,
  listing options={language=Python,basicstyle=\ttfamily\small,
    keywordstyle=\color{coulMeth}\bfseries,columns=fullflexible,
    keepspaces=true,showstringspaces=false,aboveskip=0pt,belowskip=0pt}}

\begin{remarquebox}[Durée du chapitre]
À titre indicatif, ce chapitre est prévu pour un volume de $12$~h d'enseignement
($6$ séances de $2$~h). Une partie importante des exercices est à faire en autonomie.
\end{remarquebox}

\partiecours{Limite finie, suite convergente}

Au chapitre~1, on a conjecturé sur des graphiques qu'une suite « se rapproche » d'un
nombre. On donne ici un sens précis à cette idée : la \textbf{limite} d'une suite décrit
ce que deviennent ses termes $u_{n}$ quand $n$ devient de plus en plus grand.

\smallskip
Exemple : soit la suite $(u_{n})$ définie pour tout $n\in\Ns$ par $u_{n}=2+\frac{1}{n}$.
On a $u_{1}=3$ ; $u_{10}=2{,}1$ ; $u_{100}=2{,}01$ et $u_{1\,000}=2{,}001$. L'écriture
décimale de $u_{n}$ commence par $2{,}00\ldots0$ avec de plus en plus de zéros : les termes
finissent par être aussi proches de $2$ qu'on le veut. On dit que $(u_{n})$
\textbf{converge vers $2$}.

\begin{definitionbox}[Définition 1 --- Suite convergente]
\begin{minipage}[c]{0.55\linewidth}\raggedright
\textbf{Intuitivement :} la suite $(u_{n})$ converge vers le nombre réel $\ell$ si ses termes
finissent par rester aussi près de $\ell$ qu'on le veut.

\smallskip
\textbf{Définition :} la suite $(u_{n})$ \textbf{converge} vers le nombre réel $\ell$ si
tout intervalle ouvert contenant $\ell$ contient tous les termes $u_{n}$ à partir d'un
certain rang. On note $\lim\limits_{n\to+\infty}u_{n}=\ell$ et on dit que $(u_{n})$ est
\textbf{convergente}.

\smallskip
Si elle existe, la limite est unique.
\end{minipage}\hfill
\begin{minipage}[c]{0.42\linewidth}\centering
\begin{tikzpicture}[x=0.4cm,y=0.75cm]
  \fill[red!12] (0,1.65) rectangle (12.5,2.35);
  \draw[-{Stealth[length=2mm]},thick] (0,0)--(13,0) node[below left]{$n$};
  \draw[-{Stealth[length=2mm]},thick] (0,0)--(0,3.6) node[below left]{$u_{n}$};
  \draw[dashed,black!60] (0,2)--(12.5,2) node[right,font=\footnotesize]{$\ell$};
  \draw[dotted,black!60] (4,0)--(4,3.4);
  \draw (4,0.07)--(4,-0.07) node[below,font=\footnotesize]{$N$};
  \foreach \n in {1,...,12} {%
    \pgfmathsetmacro\yy{2+2*cos(180*\n)*0.6^\n}
    \path[coulCourbe] (\n,\yy) node[croix]{};}
\end{tikzpicture}

{\footnotesize À partir du rang $N$, tous les termes sont dans la bande rouge, aussi
étroite soit-elle.}
\end{minipage}
\end{definitionbox}

\begin{methodebox}[Méthode 1 --- Conjecturer une limite et trouver un rang]
\etapeexemples
Soit la suite $(u_{n})$ définie pour tout $n\in\Ns$ par $u_{n}=2+\frac{5}{n}$.
\begin{questions}[label=\alph*)]
  \item Calculer $u_{1}$ ; $u_{10}$ ; $u_{100}$ et $u_{1\,000}$, puis conjecturer la limite
        de $(u_{n})$.
  \item Déterminer à partir de quel rang tous les termes $u_{n}$ appartiennent à
        l'intervalle $\intoo{1{,}99}{2{,}01}$.
\end{questions}
\begin{correction}
\cas{a)}
\begin{calculs}
u_{1} &= 2+5 = 7 &\qquad& \rem{on remplace $n$ par $1$}\\
u_{10} &= 2+0{,}5 = 2{,}5 &&\\
u_{100} &= 2+0{,}05 = 2{,}05 &&\\
u_{1\,000} &= 2+0{,}005 = 2{,}005 &&
\end{calculs}
\explique{Les termes se rapprochent de $2$ : on conjecture que
\res{$\lim\limits_{n\to+\infty}u_{n}=2$}.}{une conjecture, pas une preuve}
\cas{b)}
\explique{Soit $n\in\Ns$. On a $\frac{5}{n}>0$, donc $u_{n}>2$ : tous les termes sont
au-dessus de $1{,}99$.}{il reste à rester sous $2{,}01$}
\begin{calculs}
u_{n} &< 2{,}01 &\qquad& \rem{la condition cherchée}\\
2+\frac{5}{n} &< 2{,}01 &\qquad& \rem{on remplace $u_{n}$}\\
\iff \frac{5}{n} &< 0{,}01 &\qquad& \rem{on retire $2$}\\
5 &< 0{,}01\,n &\qquad& \rem{on multiplie par $n>0$}\\
\iff n &> 500 &\qquad& \rem{on divise par $0{,}01$}
\end{calculs}
\explique{\res{À partir du rang $501$}, tous les termes $u_{n}$ appartiennent à
l'intervalle $\intoo{1{,}99}{2{,}01}$.}{avec un intervalle plus étroit, on trouverait un
rang plus grand}
\end{correction}

\etape{À vous de jouer}
Soit la suite $(v_{n})$ définie pour tout $n\in\Ns$ par $v_{n}=4-\frac{3}{n}$.
\begin{questions}[label=\alph*)]
  \item Calculer $v_{1}$ ; $v_{10}$ ; $v_{100}$ et $v_{1\,000}$, puis conjecturer la limite
        de $(v_{n})$.
  \item Déterminer à partir de quel rang tous les termes $v_{n}$ appartiennent à
        l'intervalle $\intoo{3{,}99}{4{,}01}$.
\end{questions}
\reponse{a) $1$ ; $3{,}7$ ; $3{,}97$ et $3{,}997$ ; limite $4$.\; b)~À partir du rang $301$.}
\end{methodebox}

\afaire{exercices 1 à 4}

\partiecours{Limite infinie, suite divergente}

\begin{definitionbox}[Définition 2 --- Limite infinie]
\begin{minipage}[c]{0.55\linewidth}\raggedright
\textbf{Intuitivement :} la suite $(u_{n})$ tend vers $+\infty$ si ses termes finissent par
dépasser n'importe quel seuil, aussi grand soit-il, et ne plus redescendre en dessous.

\smallskip
\textbf{Définition :} la suite $(u_{n})$ \textbf{tend vers $+\infty$} si, pour tout
nombre réel $A$, l'intervalle $\intfo{A}{+\infty}$ contient tous les termes $u_{n}$ à partir d'un
certain rang. On note $\lim\limits_{n\to+\infty}u_{n}=+\infty$.

\smallskip
De même, $(u_{n})$ tend vers $-\infty$ si, pour tout nombre réel $A$, l'intervalle
$\intof{-\infty}{A}$ contient tous les termes $u_{n}$ à partir d'un certain rang.
\end{minipage}\hfill
\begin{minipage}[c]{0.42\linewidth}\centering
\begin{tikzpicture}[x=0.42cm,y=0.17cm]
  \fill[red!12] (0,10) rectangle (10.5,22);
  \draw[-{Stealth[length=2mm]},thick] (0,0)--(11,0) node[below left]{$n$};
  \draw[-{Stealth[length=2mm]},thick] (0,0)--(0,23) node[below left]{$u_{n}$};
  \draw[dashed,black!60] (0,10)--(10.5,10);
  \node[left,font=\footnotesize] at (0,10) {$A$};
  \draw[dotted,black!60] (7,0)--(7,22);
  \draw (7,0.4)--(7,-0.4) node[below,font=\footnotesize]{$N$};
  \foreach \n in {0,...,9} {%
    \pgfmathsetmacro\yy{\n*\n/4}
    \path[coulCourbe] (\n,\yy) node[croix]{};}
\end{tikzpicture}

{\footnotesize À partir du rang $N$, tous les termes dépassent $A$, aussi grand
soit-il.}
\end{minipage}
\end{definitionbox}

\begin{definitionbox}[Définition 3 --- Suite divergente]
Une suite qui n'est pas convergente est dite \textbf{divergente}. Elle peut tendre vers
$+\infty$ ou vers $-\infty$, ou n'avoir aucune limite : la suite $(w_{n})$ définie pour
tout $n\in\N$ par $w_{n}=(-1)^{n}$ vaut alternativement $1$ et $-1$ ; elle ne se
rapproche d'aucun nombre et ne dépasse aucun seuil : elle diverge sans limite.
\end{definitionbox}

Une suite qui tend vers $+\infty$ dépasse donc tout seuil $A$ : chercher le premier rang
où cela arrive est un \textbf{problème de seuil}. On le résout par le calcul quand
l'inéquation est simple, sinon avec une boucle « tant que » (vue en 1\textsuperscript{re}).

\begin{methodebox}[Méthode 2 --- Résoudre un problème de seuil]
\etapeexemples
\begin{questions}[label=\alph*)]
  \item Soit la suite $(u_{n})$ définie pour tout $n\in\N$ par $u_{n}=4n-3$. Déterminer le
        plus petit entier $n$ tel que $u_{n}\geq 10\,000$.
  \item Soit la suite $(v_{n})$ définie par $v_{0}=1$ et, pour tout $n\in\N$,
        $v_{n+1}=2v_{n}+1$ ; on admet qu'elle est croissante et tend vers $+\infty$.
        Compléter la fonction \texttt{seuil} pour qu'elle renvoie le plus petit entier $n$
        tel que $v_{n}\geq A$, puis donner la valeur renvoyée par \texttt{seuil(1000)}.
\end{questions}
\begin{python}
def seuil(A):
    n = 0
    v = 1
    while .......... :
        n = n + 1
        v = ..........
    return n
\end{python}
\begin{correction}
\cas{a)}
\begin{calculs}
4n-3 &\geq 10\,000 &\qquad& \rem{l'inéquation à résoudre}\\
\iff 4n &\geq 10\,003 &\qquad& \rem{on ajoute $3$}\\
\iff n &\geq 2\,500{,}75 &\qquad& \rem{on divise par $4>0$}
\end{calculs}
\explique{$n$ est un entier : le plus petit qui convient est \res{$n=2\,501$}.}{$2\,500$
ne convient pas}
\cas{b)}\par\nopagebreak
\begin{minipage}[t]{0.42\linewidth}\vspace{0pt}
\begin{python}
def seuil(A):
    n = 0
    v = 1
    while v < A :
        n = n + 1
        v = 2*v + 1
    return n
\end{python}
\end{minipage}\hfill
\begin{minipage}[t]{0.54\linewidth}\vspace{0pt}\raggedright
{\small\itshape\color{black!60}tant que le seuil n'est pas atteint ($v<A$), on passe au
rang suivant et on calcule le terme suivant ; la boucle s'arrête au premier rang où
$v_{n}\geq A$}
\end{minipage}
\begin{calculs}
v_{8} &= 511 &\qquad& \rem{$1$ ; $3$ ; $7$ ; $15$ ; $31$ ; $63$ ; $127$ ; $255$ ; $511$}\\
v_{9} &= 2\times 511+1 = 1\,023 &\qquad& \rem{premier terme au-delà de $1\,000$}
\end{calculs}
\explique{\texttt{seuil(1000)} renvoie \res{$9$}.}{$v_{8}<1\,000\leq v_{9}$}
\end{correction}

\etape{À vous de jouer}
\begin{questions}[label=\alph*)]
  \item Soit la suite $(a_{n})$ définie pour tout $n\in\N$ par $a_{n}=3n+7$. Déterminer le
        plus petit entier $n$ tel que $a_{n}\geq 5\,000$.
  \item Soit la suite $(b_{n})$ définie par $b_{0}=2$ et, pour tout $n\in\N$,
        $b_{n+1}=3b_{n}-1$ ; on admet qu'elle est croissante et tend vers $+\infty$.
        Écrire une fonction \texttt{seuil} qui renvoie le plus petit entier $n$ tel que
        $b_{n}\geq A$, puis donner la valeur renvoyée par \texttt{seuil(1000)}.
\end{questions}
\reponse{a) $n=1\,665$.\; b)~\texttt{while b < A}, \texttt{b = 3*b - 1} ; \texttt{seuil(1000)}
renvoie $6$.}
\end{methodebox}

\afaire{exercices 5 à 8}

\partiecours{Limites usuelles et opérations sur les limites}

\begin{proprietebox}[Propriété 1 --- Limites usuelles]
Pour tout entier $k\geq 1$ :
\begin{center}
$\lim\limits_{n\to+\infty}n^{k}=+\infty$ \qquad
$\lim\limits_{n\to+\infty}\sqrt{n}=+\infty$ \qquad
$\lim\limits_{n\to+\infty}\dfrac{1}{n^{k}}=0$ \qquad
$\lim\limits_{n\to+\infty}\dfrac{1}{\sqrt{n}}=0$
\end{center}
En particulier, $n$ et $n^{2}$ tendent vers $+\infty$ ; $\frac{1}{n}$ et $\frac{1}{n^{2}}$
tendent vers $0$. Une suite constante égale à $c$ converge vers $c$.
\end{proprietebox}

Pour une suite construite à partir de ces suites usuelles (somme, produit, quotient), on
lit sa limite dans les tableaux ci-dessous ; $\ell$ et $\ell'$ sont des nombres réels,
« $\infty$ » désigne $+\infty$ ou $-\infty$, et le signe du résultat suit la règle des
signes. Les cases « F.I. » sont des \textbf{formes indéterminées} : les tableaux ne
permettent pas de conclure (partie~IV).

\begin{proprietebox}[Propriété 2 --- Opérations sur les limites]
\textbf{Somme}

\smallskip
\begin{center}
{\renewcommand{\arraystretch}{1.8}\setlength{\tabcolsep}{7pt}%
\begin{tabular}{|c|c|c|c|c|c|c|}
\hline
$\lim u_{n}$ & $\ell$ & $\ell$ & $\ell$ & $+\infty$ & $-\infty$ & $+\infty$\\ \hline
$\lim v_{n}$ & $\ell'$ & $+\infty$ & $-\infty$ & $+\infty$ & $-\infty$ & $-\infty$\\ \hline
$\lim (u_{n}+v_{n})$ & $\ell+\ell'$ & $+\infty$ & $-\infty$ & $+\infty$ & $-\infty$ &
\textbf{F.I.}\\ \hline
\end{tabular}}
\end{center}

\textbf{Produit}

\smallskip
\begin{center}
{\renewcommand{\arraystretch}{1.8}\setlength{\tabcolsep}{7pt}%
\begin{tabular}{|c|c|c|c|c|}
\hline
$\lim u_{n}$ & $\ell$ & $\ell\neq 0$ & $\infty$ & $0$\\ \hline
$\lim v_{n}$ & $\ell'$ & $\infty$ & $\infty$ & $\infty$\\ \hline
$\lim (u_{n}\times v_{n})$ & $\ell\times\ell'$ & $\infty$ & $\infty$ & \textbf{F.I.}\\ \hline
\end{tabular}}
\end{center}

\textbf{Quotient}

\smallskip
\begin{center}
{\renewcommand{\arraystretch}{2.4}\setlength{\tabcolsep}{6pt}%
\begin{tabular}{|c|c|c|c|c|c|c|}
\hline
$\lim u_{n}$ & $\ell$ & $\ell$ & $\ell\neq 0$ ou $\infty$ & $\infty$ & $\infty$ & $0$\\ \hline
$\lim v_{n}$ & $\ell'\neq 0$ & $\infty$ & $0$, $v_{n}$ de signe constant & $\ell'$ & $\infty$ & $0$\\ \hline
$\lim \dfrac{u_{n}}{v_{n}}$ & $\dfrac{\ell}{\ell'}$ & $0$ & $\infty$ & $\infty$ &
\textbf{F.I.} & \textbf{F.I.}\\ \hline
\end{tabular}}
\end{center}

\textbf{En mots :} un nombre divisé par « très grand » devient « très petit » ; un nombre
non nul divisé par « très proche de zéro » devient « très grand ».
\end{proprietebox}

\begin{methodebox}[Méthode 3 --- Calculer une limite par opérations]
\etapeexemples
Déterminer la limite de chacune des suites $(u_{n})$, $(v_{n})$, $(w_{n})$ définies pour
tout $n\in\Ns$ par :
\begin{center}
a)~$u_{n}=n^{2}+3\sqrt{n}-1$ \qquad
b)~$v_{n}=\left(2+\dfrac{1}{n}\right)\left(5-n^{2}\right)$ \qquad
c)~$w_{n}=\dfrac{-3}{n^{2}+1}$
\end{center}
\begin{correction}
\cas{a)}
\explique{$\lim\limits_{n\to+\infty}\sqrt{n}=+\infty$, donc
$\lim\limits_{n\to+\infty}3\sqrt{n}=+\infty$.}{produit par $3>0$}
\explique{$\lim\limits_{n\to+\infty}n^{2}=+\infty$, donc par somme,
\res{$\lim\limits_{n\to+\infty}u_{n}=+\infty$}.}{$+\infty+\infty-1$}
\cas{b)}
\explique{$\lim\limits_{n\to+\infty}\frac{1}{n}=0$, donc par somme,
$\lim\limits_{n\to+\infty}\left(2+\frac{1}{n}\right)=2$.}{}
\explique{$\lim\limits_{n\to+\infty}n^{2}=+\infty$, donc par somme,
$\lim\limits_{n\to+\infty}\left(5-n^{2}\right)=-\infty$.}{}
\explique{Par produit, \res{$\lim\limits_{n\to+\infty}v_{n}=-\infty$}.}{$2\times(-\infty)$ :
règle des signes}
\cas{c)}
\explique{$\lim\limits_{n\to+\infty}n^{2}=+\infty$, donc par somme,
$\lim\limits_{n\to+\infty}\left(n^{2}+1\right)=+\infty$.}{}
\explique{Par quotient, \res{$\lim\limits_{n\to+\infty}w_{n}=0$}.}{$-3$ divisé par
« très grand »}
\end{correction}

\etape{À vous de jouer}
Déterminer la limite de chacune des suites $(a_{n})$, $(b_{n})$, $(c_{n})$ définies pour
tout $n\geq 3$ par :
\begin{center}
a)~$a_{n}=2n+\sqrt{n}-4$ \qquad
b)~$b_{n}=\left(3-\dfrac{1}{n^{2}}\right)\left(n+4\right)$ \qquad
c)~$c_{n}=\dfrac{5}{2-n}$
\end{center}
\reponse{a) $+\infty$.\; b)~$+\infty$.\; c)~$0$.}
\end{methodebox}

\afaire{exercices 9 à 12}

\partiecours{Lever une indétermination}

Les quatre formes indéterminées sont « $\infty-\infty$ », « $0\times\infty$ »,
« $\frac{\infty}{\infty}$ » et « $\frac{0}{0}$ ». Une forme indéterminée ne veut pas dire
qu'il n'y a pas de limite : il faut transformer l'écriture de $u_{n}$ pour sortir de
l'impasse, puis appliquer les opérations. Pour une somme de puissances de $n$, on
\textbf{factorise par le terme de plus haut degré}, celui qui « grandit le plus vite »
($n^{2}$ dans $2n^{2}-7n+3$).

\begin{methodebox}[Méthode 4 --- Factoriser par le terme de plus haut degré]
\etapeexemples
Déterminer la limite de chacune des suites $(u_{n})$, $(v_{n})$, $(w_{n})$ définies pour
tout $n\in\Ns$ par :
\begin{center}
a)~$u_{n}=2n^{2}-7n+3$ \qquad
b)~$v_{n}=\dfrac{6n^{2}+1}{3n^{2}+2n}$ \qquad
c)~$w_{n}=n-4\sqrt{n}$
\end{center}
\begin{correction}
\cas{a)}
\explique{$\lim\limits_{n\to+\infty}2n^{2}=+\infty$ et
$\lim\limits_{n\to+\infty}(-7n+3)=-\infty$ : c'est une forme indéterminée
« $\infty-\infty$ ».}{on le signale d'abord}
\explique{Pour tout $n\in\Ns$, on a $n\neq 0$ :}{on peut factoriser par $n^{2}$}
\begin{calculs}
u_{n} &= n^{2}\left(2-\frac{7n}{n^{2}}+\frac{3}{n^{2}}\right) &\qquad& \rem{on factorise
par $n^{2}$, terme de plus haut degré}\\
u_{n} &= n^{2}\left(2-\frac{7}{n}+\frac{3}{n^{2}}\right) &\qquad& \rem{on simplifie
$\frac{7n}{n^{2}}$}
\end{calculs}
\explique{$\lim\limits_{n\to+\infty}\frac{7}{n}=0$ et $\lim\limits_{n\to+\infty}\frac{3}{n^{2}}=0$,
donc par somme, $\lim\limits_{n\to+\infty}\left(2-\frac{7}{n}+\frac{3}{n^{2}}\right)=2$.}{}
\explique{$\lim\limits_{n\to+\infty}n^{2}=+\infty$, donc par produit,
\res{$\lim\limits_{n\to+\infty}u_{n}=+\infty$}.}{$(+\infty)\times 2$}
\cas{b)}
\explique{Le numérateur et le dénominateur tendent vers $+\infty$ : c'est une forme
indéterminée « $\frac{\infty}{\infty}$ ».}{}
\explique{Pour tout $n\in\Ns$, on a $n\neq 0$ :}{on peut factoriser et simplifier par $n^{2}$}
\begin{calculs}
v_{n} &= \frac{n^{2}\left(6+\frac{1}{n^{2}}\right)}{n^{2}\left(3+\frac{2}{n}\right)} &\qquad&
\rem{on factorise en haut et en bas par $n^{2}$}\\
v_{n} &= \frac{6+\frac{1}{n^{2}}}{3+\frac{2}{n}} &\qquad& \rem{on simplifie par
$n^{2}\neq 0$}
\end{calculs}
\explique{Par somme, le numérateur tend vers $6$ et le dénominateur vers $3$ ; par
quotient, \res{$\lim\limits_{n\to+\infty}v_{n}=2$}.}{$\frac{6}{3}=2$}
\cas{c)}
\explique{$\lim\limits_{n\to+\infty}n=+\infty$ et
$\lim\limits_{n\to+\infty}4\sqrt{n}=+\infty$ : forme indéterminée
« $\infty-\infty$ ».}{}
\explique{Pour tout $n\in\Ns$, on a $n\neq 0$ :}{on peut factoriser par $n$}
\begin{calculs}
w_{n} &= n\left(1-\frac{4\sqrt{n}}{n}\right) &\qquad& \rem{on factorise par $n$}\\
w_{n} &= n\left(1-\frac{4}{\sqrt{n}}\right) &\qquad& \rem{$n=\sqrt{n}\times\sqrt{n}$}
\end{calculs}
\explique{$\lim\limits_{n\to+\infty}\frac{4}{\sqrt{n}}=0$, donc par somme, la parenthèse tend
vers $1$ ; par produit, \res{$\lim\limits_{n\to+\infty}w_{n}=+\infty$}.}{$(+\infty)\times 1$}
\end{correction}

\etape{À vous de jouer}
Déterminer la limite de chacune des suites $(a_{n})$, $(b_{n})$, $(c_{n})$ définies pour
tout $n\in\Ns$ par :
\begin{center}
a)~$a_{n}=-3n^{2}+4n+7$ \qquad
b)~$b_{n}=\dfrac{5n^{2}-1}{2n^{2}+n}$ \qquad
c)~$c_{n}=2n-3\sqrt{n}$
\end{center}
\reponse{a) $-\infty$.\; b)~$\frac{5}{2}$.\; c)~$+\infty$.}
\end{methodebox}

\begin{remarquebox}[Quotient de degrés différents]
La même factorisation conclut : pour tout $n\in\Ns$ ($n\neq 0$), on a
$\frac{n^{2}+4}{2n+1}=n\times\frac{1+\frac{4}{n^{2}}}{2+\frac{1}{n}}$, qui tend vers
$+\infty$ par produit ; et $\frac{3n+1}{n^{2}+2}=\frac{1}{n}\times\frac{3+\frac{1}{n}}{1+\frac{2}{n^{2}}}$,
qui tend vers $0$. Le terme de plus haut degré du numérateur et celui du dénominateur décident.
\end{remarquebox}

Avec des racines carrées, la forme « $\infty-\infty$ » de $\sqrt{n+5}-\sqrt{n}$ ne se lève
pas en factorisant. On utilise l'identité remarquable $(a-b)(a+b)=a^{2}-b^{2}$, qui fait
disparaître les racines : $\left(\sqrt{n+5}-\sqrt{n}\right)\left(\sqrt{n+5}+\sqrt{n}\right)=(n+5)-n$.
L'expression $\sqrt{n+5}+\sqrt{n}$ est la \textbf{quantité conjuguée} de
$\sqrt{n+5}-\sqrt{n}$.

\begin{methodebox}[Méthode 5 --- Multiplier par la quantité conjuguée]
\etapeexemples
Soit la suite $(u_{n})$ définie pour tout $n\in\N$ par $u_{n}=\sqrt{n+5}-\sqrt{n}$.
Déterminer sa limite.
\begin{correction}
\cas{Résolution}\par\nopagebreak
\explique{$\lim\limits_{n\to+\infty}\sqrt{n+5}=+\infty$ et
$\lim\limits_{n\to+\infty}\sqrt{n}=+\infty$ : forme indéterminée « $\infty-\infty$ ».}{}
\begin{calculs}
u_{n} &= \frac{\left(\sqrt{n+5}-\sqrt{n}\right)\left(\sqrt{n+5}+\sqrt{n}\right)}{\sqrt{n+5}+\sqrt{n}}
&\qquad& \rem{on multiplie et on divise par la quantité conjuguée}\\
u_{n} &= \frac{(n+5)-n}{\sqrt{n+5}+\sqrt{n}} &\qquad& \rem{$(a-b)(a+b)=a^{2}-b^{2}$}\\
u_{n} &= \frac{5}{\sqrt{n+5}+\sqrt{n}} &&
\end{calculs}
\explique{Par somme, $\lim\limits_{n\to+\infty}\left(\sqrt{n+5}+\sqrt{n}\right)=+\infty$ ;
par quotient, \res{$\lim\limits_{n\to+\infty}u_{n}=0$}.}{$5$ divisé par « très grand »}
\end{correction}

\etape{À vous de jouer}
Soit la suite $(v_{n})$ définie pour tout $n\in\N$ par $v_{n}=\sqrt{n+3}-\sqrt{n+1}$.
Déterminer sa limite.
\reponse{$v_{n}=\frac{2}{\sqrt{n+3}+\sqrt{n+1}}$ ; limite $0$.}
\end{methodebox}

\afaire{exercices 13 à 17}

\partiecours{Limites et comparaison}

Les termes d'une suite comme $n^{2}-3\cos(n)$ ne se calculent pas avec les tableaux :
$\cos(n)$ n'a pas de limite. Mais $\cos(n)$ reste entre $-1$ et $1$ : on compare la
suite à une suite plus simple.

\begin{proprietebox}[Théorème 1 --- Théorèmes de comparaison]
\begin{minipage}[c]{0.58\linewidth}\raggedright
Soit $(u_{n})$ et $(v_{n})$ deux suites telles que, à partir d'un certain rang,
$u_{n}\leq v_{n}$.
\begin{itemize}
  \item Si $\lim\limits_{n\to+\infty}u_{n}=+\infty$, alors
        $\lim\limits_{n\to+\infty}v_{n}=+\infty$ : $(u_{n})$ « pousse » $(v_{n})$ vers
        le haut.
  \item Si $\lim\limits_{n\to+\infty}v_{n}=-\infty$, alors
        $\lim\limits_{n\to+\infty}u_{n}=-\infty$ : $(v_{n})$ « pousse » $(u_{n})$ vers
        le bas.
\end{itemize}
\end{minipage}\hfill
\begin{minipage}[c]{0.38\linewidth}\centering
\begin{tikzpicture}[x=0.4cm,y=0.45cm]
  \draw[-{Stealth[length=2mm]},thick] (0,0)--(12,0) node[below left]{$n$};
  \draw[-{Stealth[length=2mm]},thick] (0,0)--(0,7.5);
  \foreach \n in {0,...,11} {%
    \pgfmathsetmacro\yy{0.5*\n}
    \pgfmathsetmacro\zz{0.5*\n+1+0.6*cos(180*\n)}
    \path[coulCourbe] (\n,\yy) node[croix]{};
    \path[coulCourbeB] (\n,\zz) node[croixdroite]{};}
  \node[coulCourbe,font=\footnotesize] at (10,3.8) {$u_{n}$};
  \node[coulCourbeB,font=\footnotesize] at (8.2,6.9) {$v_{n}$};
\end{tikzpicture}
\end{minipage}

\smallskip
\textbf{\color{coulMeth}Démonstration} (exigible) du premier point.
Soit $A$ un nombre réel. Comme $\lim\limits_{n\to+\infty}u_{n}=+\infty$, il existe un rang
$n_{1}$ à partir duquel $u_{n}\geq A$. Par hypothèse, il existe un rang $n_{2}$ à partir
duquel $u_{n}\leq v_{n}$. Pour $n$ supérieur au plus grand des deux rangs $n_{1}$ et
$n_{2}$, les deux inégalités sont vraies :
\begin{calculs}
v_{n} &\geq u_{n}\geq A &\qquad& \rem{à partir de ce rang}
\end{calculs}
L'intervalle $\intfo{A}{+\infty}$ contient donc tous les termes $v_{n}$ à partir d'un
certain rang, et cela pour tout nombre réel $A$ : $\lim\limits_{n\to+\infty}v_{n}=+\infty$.

\end{proprietebox}

\begin{methodebox}[Méthode 6 --- Démontrer une limite infinie par comparaison]
\etapeexemples
Déterminer la limite de chacune des suites $(u_{n})$ et $(v_{n})$ définies pour tout
$n\in\N$ par :
\begin{center}
a)~$u_{n}=n^{2}-3\cos(n)$ \qquad b)~$v_{n}=-n+\sin(n)$
\end{center}
\begin{correction}
\cas{a)}
\explique{Soit $n\in\N$.}{}
\begin{calculs}
\cos(n) &\leq 1 &\qquad& \rem{un cosinus est entre $-1$ et $1$}\\
\iff -3\cos(n) &\geq -3 &\qquad& \rem{on multiplie par $-3<0$ : l'ordre change}\\
\iff n^{2}-3\cos(n) &\geq n^{2}-3 &\qquad& \rem{on ajoute $n^{2}$}\\
u_{n} &\geq n^{2}-3 &\qquad& \rem{on reconnaît $u_{n}$}
\end{calculs}
\explique{$\lim\limits_{n\to+\infty}\left(n^{2}-3\right)=+\infty$, donc par comparaison,
\res{$\lim\limits_{n\to+\infty}u_{n}=+\infty$}.}{$(u_{n})$ est au-dessus d'une suite qui
tend vers $+\infty$}
\cas{b)}
\explique{Soit $n\in\N$.}{}
\begin{calculs}
\sin(n) &\leq 1 &\qquad& \rem{un sinus est entre $-1$ et $1$}\\
\iff -n+\sin(n) &\leq -n+1 &\qquad& \rem{on ajoute $-n$}\\
v_{n} &\leq -n+1 &\qquad& \rem{on reconnaît $v_{n}$}
\end{calculs}
\explique{$\lim\limits_{n\to+\infty}(-n+1)=-\infty$, donc par comparaison,
\res{$\lim\limits_{n\to+\infty}v_{n}=-\infty$}.}{$(v_{n})$ est en dessous d'une suite
qui tend vers $-\infty$}
\end{correction}

\etape{À vous de jouer}
Déterminer la limite de chacune des suites $(a_{n})$ et $(b_{n})$ définies pour tout
$n\in\N$ par :
\begin{center}
a)~$a_{n}=3n+2\sin(n)$ \qquad b)~$b_{n}=5-n^{2}+(-1)^{n}$
\end{center}
\reponse{a) $a_{n}\geq 3n-2$ : $+\infty$.\; b)~$b_{n}\leq 6-n^{2}$ : $-\infty$.}
\end{methodebox}

Pour une limite finie, une seule suite ne suffit pas : il faut « coincer » la suite entre
deux suites qui ont la même limite.

\begin{proprietebox}[Théorème 2 --- Théorème des gendarmes (admis)]
\begin{minipage}[c]{0.58\linewidth}\raggedright
Soit $(u_{n})$, $(v_{n})$ et $(w_{n})$ trois suites telles que, à partir d'un certain rang,
$u_{n}\leq v_{n}\leq w_{n}$. Si $(u_{n})$ et $(w_{n})$ convergent vers le même nombre réel
$\ell$, alors $(v_{n})$ converge aussi vers $\ell$.

\smallskip
Les deux « gendarmes » $(u_{n})$ et $(w_{n})$ se resserrent autour de $\ell$ et emmènent
$(v_{n})$ avec eux.
\end{minipage}\hfill
\begin{minipage}[c]{0.38\linewidth}\centering
\begin{tikzpicture}[x=0.42cm,y=0.85cm]
  \draw[-{Stealth[length=2mm]},thick] (0,0)--(12,0) node[below left]{$n$};
  \draw[-{Stealth[length=2mm]},thick] (0,0)--(0,4.6);
  \draw[dashed,black!60] (0,2)--(11.5,2) node[right,font=\footnotesize]{$\ell$};
  \foreach \n in {0,...,11} {%
    \pgfmathsetmacro\aa{2-2/(\n+1)}
    \pgfmathsetmacro\bb{2+2/(\n+1)}
    \pgfmathsetmacro\cc{2+1.5*cos(180*\n)/(\n+1)}
    \path[coulCourbeB] (\n,\aa) node[croixdroite]{} (\n,\bb) node[croixdroite]{};
    \path[coulCourbe] (\n,\cc) node[croix]{};}
\end{tikzpicture}
\end{minipage}
\end{proprietebox}

\begin{methodebox}[Méthode 7 --- Démontrer une limite finie par encadrement]
\etapeexemples
Soit la suite $(u_{n})$ définie pour tout $n\in\Ns$ par $u_{n}=3+\dfrac{2\cos(n)}{n}$.
Déterminer sa limite.
\begin{correction}
\cas{Résolution}\par\nopagebreak
\explique{Soit $n\in\Ns$.}{}
\begin{calculs}
-1 &\leq \cos(n)\leq 1 &\qquad& \rem{on part de l'encadrement du cosinus}\\
-2 &\leq 2\cos(n)\leq 2 &\qquad& \rem{on multiplie par $2>0$}\\
\iff -\frac{2}{n} &\leq \frac{2\cos(n)}{n}\leq \frac{2}{n} &\qquad& \rem{on divise par $n>0$}\\
\iff 3-\frac{2}{n} &\leq 3+\frac{2\cos(n)}{n}\leq 3+\frac{2}{n} &\qquad& \rem{on ajoute $3$}\\
3-\frac{2}{n} &\leq u_{n}\leq 3+\frac{2}{n} &\qquad& \rem{on reconnaît $u_{n}$}
\end{calculs}
\explique{$\lim\limits_{n\to+\infty}\frac{2}{n}=0$, donc par somme, les suites
$\left(3-\frac{2}{n}\right)$ et $\left(3+\frac{2}{n}\right)$ convergent vers $3$.}{les deux
gendarmes ont la même limite}
\explique{D'après le théorème des gendarmes, \res{$\lim\limits_{n\to+\infty}u_{n}=3$}.}{}
\end{correction}

\etape{À vous de jouer}
Soit la suite $(v_{n})$ définie pour tout $n\in\Ns$ par $v_{n}=1-\dfrac{\sin(n)}{n^{2}}$.
Déterminer sa limite.
\reponse{$1-\frac{1}{n^{2}}\leq v_{n}\leq 1+\frac{1}{n^{2}}$ ; limite $1$.}
\end{methodebox}

\begin{proprietebox}[Propriété 3 --- Passage à la limite dans une inégalité (admis)]
Soit $(u_{n})$ et $(v_{n})$ deux suites convergentes, de limites $\ell$ et $\ell'$. Si, à
partir d'un certain rang, $u_{n}\leq v_{n}$, alors $\ell\leq\ell'$. En particulier, si
pour tout $n\in\N$, $u_{n}\leq M$, alors $\ell\leq M$.

\smallskip
Attention : une inégalité stricte devient large. Pour tout $n\in\Ns$, $\frac{1}{n}>0$,
mais la limite vaut $0$.
\end{proprietebox}

\afaire{exercices 18 à 22}

\partiecours{Convergence des suites monotones}

Au chapitre~1, on a vu comment démontrer par récurrence qu'une suite est croissante et
majorée (définition~4 : $(u_{n})$ est majorée par $M$ si, pour tout $n\in\N$,
$u_{n}\leq M$). Ici, on ajoute : une suite qui monte sans jamais dépasser un plafond
finit par se stabiliser ; elle converge.

\begin{proprietebox}[Théorème 3 --- Théorème de convergence monotone (admis)]
\begin{itemize}
  \item Toute suite croissante et majorée converge.
  \item Toute suite décroissante et minorée converge.
\end{itemize}
Le théorème dit que la suite converge, sans donner sa limite : une suite croissante
majorée par $3$ converge vers un nombre $\ell\leq 3$ (propriété~3), mais pas
forcément vers $3$.
\end{proprietebox}

\begin{proprietebox}[Propriété 4 --- Suite croissante non majorée]
Toute suite croissante et non majorée tend vers $+\infty$. (De même, toute suite
décroissante et non minorée tend vers $-\infty$.)

\smallskip
\textbf{\color{coulMeth}Démonstration} (exigible).
Soit $(u_{n})$ une suite croissante non majorée, et $A$ un nombre réel. Comme $A$ n'est
pas un majorant de la suite, il existe un rang $p$ tel que $u_{p}>A$. La suite est
croissante, donc pour tout $n\geq p$ :
\begin{calculs}
u_{n} &\geq u_{p}>A &\qquad& \rem{$(u_{n})$ croissante}
\end{calculs}
L'intervalle $\intfo{A}{+\infty}$ contient tous les termes à partir du rang $p$, et cela
pour tout nombre réel $A$ : $\lim\limits_{n\to+\infty}u_{n}=+\infty$.

\end{proprietebox}

\begin{methodebox}[Méthode 8 --- Démontrer qu'une suite converge]
\etapeexemples
Soit la suite $(u_{n})$ définie par $u_{0}=0$ et, pour tout $n\in\N$,
$u_{n+1}=0{,}8\,u_{n}+4$.
\begin{questions}[label=\alph*)]
  \item Démontrer par récurrence que, pour tout $n\in\N$, on a $u_{n}\leq u_{n+1}\leq 20$.
  \item En déduire que la suite $(u_{n})$ converge, et donner une majoration de sa
        limite $\ell$.
\end{questions}
\begin{correction}
\cas{a)}
\explique{Pour tout $n\in\N$, on note $P_{n}$ la propriété :
«~$u_{n}\leq u_{n+1}\leq 20$~».}{}
\explique{\textbf{Initialisation.} $u_{0}=0$ et $u_{1}=0{,}8\times 0+4=4$ ;
$0\leq 4\leq 20$, donc $P_{0}$ est vraie.}{}
\explique{\textbf{Hérédité.} Soit $n\in\N$ tel que $P_{n}$ est vraie. Montrons que
$u_{n+1}\leq u_{n+2}\leq 20$.}{}
\begin{calculs}
u_{n} &\leq u_{n+1}\leq 20 &\qquad& \rem{hypothèse de récurrence}\\
0{,}8\,u_{n} &\leq 0{,}8\,u_{n+1}\leq 16 &\qquad& \rem{on multiplie par $0{,}8>0$}\\
0{,}8\,u_{n}+4 &\leq 0{,}8\,u_{n+1}+4\leq 20 &\qquad& \rem{on ajoute $4$}\\
u_{n+1} &\leq u_{n+2}\leq 20 &\qquad& \rem{relation de récurrence}
\end{calculs}
\explique{Donc $P_{n+1}$ est vraie : la propriété est héréditaire.}{}
\explique{\textbf{Conclusion.} $P_{0}$ est vraie et la propriété est héréditaire ;
d'après le principe de récurrence, \res{pour tout $n\in\N$, on a
$u_{n}\leq u_{n+1}\leq 20$}.}{}
\cas{b)}
\explique{Pour tout $n\in\N$, on a $u_{n}\leq u_{n+1}$ : la suite $(u_{n})$ est
croissante ; et $u_{n}\leq 20$ : elle est majorée par $20$.}{on lit les deux informations}
\explique{D'après le théorème de convergence monotone, \res{$(u_{n})$ converge} vers un
nombre réel $\ell$, et \res{$\ell\leq 20$}.}{propriété~3 ; la valeur de $\ell$ : partie~X}
\end{correction}

\etape{À vous de jouer}
Soit la suite $(v_{n})$ définie par $v_{0}=1$ et, pour tout $n\in\N$,
$v_{n+1}=0{,}5\,v_{n}+3$.
\begin{questions}[label=\alph*)]
  \item Démontrer par récurrence que, pour tout $n\in\N$, on a $v_{n}\leq v_{n+1}\leq 6$.
  \item En déduire que la suite $(v_{n})$ converge, et donner une majoration de sa
        limite $\ell$.
\end{questions}
\reponse{a) $v_{1}=3{,}5$ ; hérédité : $\times 0{,}5$ puis $+3$.\; b)~Croissante et majorée
par $6$ : elle converge, et $\ell\leq 6$.}
\end{methodebox}

\afaire{exercices 23 à 26}

\partiecours{Limite d'une suite géométrique}

Rappel de 1\textsuperscript{re} : une suite géométrique de raison $q$ et de premier terme
$u_{0}$ a pour terme général $u_{n}=u_{0}\times q^{n}$. Sa limite dépend donc de celle de
$q^{n}$ : multiplier sans fin par $2$ fait « exploser » les termes, multiplier sans fin
par $0{,}5$ les « écrase » vers $0$.

\begin{proprietebox}[Propriété 5 --- Limite de $q^{n}$]
Soit $q$ un nombre réel.
\begin{center}
{\renewcommand{\arraystretch}{1.8}\setlength{\tabcolsep}{8pt}%
\begin{tabular}{|c|c|c|c|c|}
\hline
$q$ & $q\leq -1$ & $-1<q<1$ & $q=1$ & $q>1$\\ \hline
$\lim\limits_{n\to+\infty}q^{n}$ & pas de limite & $0$ & $1$ & $+\infty$\\ \hline
\end{tabular}}
\end{center}
Exemples : $\lim\limits_{n\to+\infty}3^{n}=+\infty$ ; $\lim\limits_{n\to+\infty}0{,}9^{n}=0$ ;
$\lim\limits_{n\to+\infty}(-0{,}5)^{n}=0$ ; la suite $\left((-2)^{n}\right)$ prend des
valeurs alternativement positives et négatives de plus en plus grandes : elle n'a pas de
limite.

\smallskip
\textbf{\color{coulMeth}Démonstration} (exigible) pour $q>1$.
On pose $a=q-1$, de sorte que $a>0$ et $q=1+a$. D'après l'inégalité de Bernoulli
(chapitre~1), pour tout $n\in\N$ :
\begin{calculs}
q^{n} &= (1+a)^{n}\geq 1+na &\qquad& \rem{inégalité de Bernoulli, $a>0$}
\end{calculs}
Or $\lim\limits_{n\to+\infty}(1+na)=+\infty$ car $a>0$. Par comparaison,
$\lim\limits_{n\to+\infty}q^{n}=+\infty$.

\smallskip
Pour $0<q<1$ : $\frac{1}{q}>1$, donc $\left(\frac{1}{q}\right)^{n}$ tend vers $+\infty$, et
$q^{n}=\dfrac{1}{\left(\frac{1}{q}\right)^{n}}$ tend vers $0$ par quotient.

\end{proprietebox}

\begin{methodebox}[Méthode 9 --- Déterminer la limite d'une suite géométrique]
\etapeexemples
Déterminer la limite de chacune des suites $(u_{n})$, $(v_{n})$, $(w_{n})$ définies pour
tout $n\in\N$ par :
\begin{center}
a)~$u_{n}=5\times 1{,}2^{n}-3$ \qquad
b)~$v_{n}=4-2\times 0{,}6^{n}$ \qquad
c)~$w_{n}=2+(-0{,}5)^{n}$
\end{center}
\begin{correction}
\cas{a)}
\explique{$1{,}2>1$, donc $\lim\limits_{n\to+\infty}1{,}2^{n}=+\infty$.}{propriété~5}
\explique{Par produit par $5>0$, puis par somme, \res{$\lim\limits_{n\to+\infty}u_{n}=+\infty$}.}{}
\cas{b)}
\explique{$-1<0{,}6<1$, donc $\lim\limits_{n\to+\infty}0{,}6^{n}=0$.}{}
\explique{Par produit, $\lim\limits_{n\to+\infty}\left(-2\times 0{,}6^{n}\right)=0$ ; par
somme, \res{$\lim\limits_{n\to+\infty}v_{n}=4$}.}{$4-0$}
\cas{c)}
\explique{$-1<-0{,}5<1$, donc $\lim\limits_{n\to+\infty}(-0{,}5)^{n}=0$ ; par somme,
\res{$\lim\limits_{n\to+\infty}w_{n}=2$}.}{une raison négative entre $-1$ et $1$ donne
aussi $0$}
\end{correction}

\etape{À vous de jouer}
Déterminer la limite de chacune des suites $(a_{n})$, $(b_{n})$, $(c_{n})$ définies pour
tout $n\in\N$ par :
\begin{center}
a)~$a_{n}=3\times 1{,}5^{n}+1$ \qquad
b)~$b_{n}=7-5\times 0{,}8^{n}$ \qquad
c)~$c_{n}=1-(-0{,}4)^{n}$
\end{center}
\reponse{a) $+\infty$.\; b)~$7$.\; c)~$1$.}
\end{methodebox}

\afaire{exercices 27 à 29}

\partiecours{Sommes de termes (rappels de 1\textsuperscript{re}) et limite}

\begin{proprietebox}[Rappel 1 --- Sommes de termes consécutifs]
\begin{itemize}
  \item Pour tout $n\in\Ns$, on a $1+2+\dots+n=\dfrac{n(n+1)}{2}$. Pour une suite
        arithmétique, la somme de termes consécutifs vaut
        $\text{(nombre de termes)}\times\dfrac{\text{premier terme}+\text{dernier terme}}{2}$.
  \item Pour tout nombre réel $q\neq 1$ et tout $n\in\N$, on a
        $1+q+q^{2}+\dots+q^{n}=\dfrac{1-q^{n+1}}{1-q}$.
  \item Pour une suite géométrique $(u_{n})$ de raison $q\neq 1$ :
        $u_{0}+u_{1}+\dots+u_{n}=u_{0}\times\dfrac{1-q^{n+1}}{1-q}$ ; la somme compte
        $n+1$ termes, d'où l'exposant $n+1$.
\end{itemize}
Exemple : $1+2+\dots+100=\frac{100\times 101}{2}$, soit $5\,050$ ; et
$1+3+3^{2}+3^{3}=\frac{1-3^{4}}{1-3}$, soit $40$.
\end{proprietebox}

Quand $-1<q<1$, le terme $q^{n+1}$ tend vers $0$ : la somme d'une infinité de termes de
plus en plus petits peut rester finie.

\begin{methodebox}[Méthode 10 --- Déterminer la limite d'une somme de termes géométriques]
\etapeexemples
Soit $(u_{n})$ la suite géométrique de raison $0{,}5$ et de premier terme $u_{0}=6$. Pour
tout $n\in\N$, on note $S_{n}=u_{0}+u_{1}+\dots+u_{n}$.
\begin{questions}[label=\alph*)]
  \item Calculer $S_{3}$.
  \item Exprimer $S_{n}$ en fonction de $n$.
  \item Déterminer la limite de la suite $(S_{n})$.
\end{questions}
\begin{correction}
\cas{a)}
\begin{calculs}
S_{3} &= 6+3+1{,}5+0{,}75 &\qquad& \rem{chaque terme est la moitié du précédent}\\
S_{3} &= \res{11{,}25} &&
\end{calculs}
\cas{b)}
\begin{calculs}
S_{n} &= u_{0}\times\frac{1-q^{n+1}}{1-q} &\qquad& \rem{rappel~1}\\
S_{n} &= 6\times\frac{1-0{,}5^{n+1}}{0{,}5} &\qquad& \rem{$1-0{,}5=0{,}5$}\\
S_{n} &= \res{12\left(1-0{,}5^{n+1}\right)} &\qquad& \rem{diviser par $0{,}5$, c'est
multiplier par $2$}
\end{calculs}
\explique{Contrôle : $S_{3}=12\times(1-0{,}0625)$, soit $11{,}25$.}{on vérifie avec a)}
\cas{c)}
\explique{$-1<0{,}5<1$, donc $\lim\limits_{n\to+\infty}0{,}5^{n+1}=0$.}{$0{,}5^{n+1}=0{,}5\times 0{,}5^{n}$}
\explique{Par somme puis par produit, \res{$\lim\limits_{n\to+\infty}S_{n}=12$}.}{$12\times(1-0)$}
\end{correction}

\etape{À vous de jouer}
Soit $(v_{n})$ la suite géométrique de raison $0{,}2$ et de premier terme $v_{0}=4$. Pour
tout $n\in\N$, on note $T_{n}=v_{0}+v_{1}+\dots+v_{n}$.
\begin{questions}[label=\alph*)]
  \item Calculer $T_{2}$.
  \item Exprimer $T_{n}$ en fonction de $n$.
  \item Déterminer la limite de la suite $(T_{n})$.
\end{questions}
\reponse{a) $4{,}96$.\; b)~$T_{n}=5\left(1-0{,}2^{n+1}\right)$.\; c)~$5$.}
\end{methodebox}

\afaire{exercices 30 à 33}

\partiecours{Suites arithmético-géométriques}

\begin{definitionbox}[Définition 4 --- Suite arithmético-géométrique]
Une suite $(u_{n})$ est \textbf{arithmético-géométrique} s'il existe deux nombres réels
$a$ et $b$ tels que, pour tout $n\in\N$, $u_{n+1}=a\,u_{n}+b$.

\smallskip
Elle modélise une évolution « en pourcentage, plus une quantité fixe » : un parc qui perd
$10\,\%$ de ses vélos et en reçoit $120$ chaque mois vérifie $u_{n+1}=0{,}9\,u_{n}+120$.
\end{definitionbox}

\begin{proprietebox}[Propriété 6 --- Suite auxiliaire géométrique]
Soit $(u_{n})$ une suite telle que, pour tout $n\in\N$, $u_{n+1}=a\,u_{n}+b$, avec
$a\neq 1$. Soit $c$ le nombre réel tel que $c=a\,c+b$, c'est-à-dire
$c=\dfrac{b}{1-a}$. Alors la suite $(v_{n})$ définie pour tout $n\in\N$ par
$v_{n}=u_{n}-c$ est géométrique de raison $a$, et :
\begin{center}
$u_{n}=(u_{0}-c)\times a^{n}+c$.
\end{center}
Si $-1<a<1$, la suite $(u_{n})$ converge vers $c$.

\smallskip
En effet, $v_{n+1}=a\,u_{n}+b-c$ ; or $b-c=-a\,c$, donc $v_{n+1}=a\,(u_{n}-c)$, soit
$a\,v_{n}$. Au bac, l'énoncé donne la suite $(v_{n})$ : on refait ce calcul (chapitre~1,
méthode~3).
\end{proprietebox}

\begin{methodebox}[Méthode 11 --- Étudier une suite arithmético-géométrique]
\etapeexemples
Une ville met $800$ vélos en libre-service en janvier. Chaque mois, $10\,\%$ des vélos
partent en révision et ne reviennent pas, et la ville en ajoute $120$ neufs. On note
$u_{n}$ le nombre de vélos $n$ mois après janvier : $u_{0}=800$ et, pour tout $n\in\N$,
$u_{n+1}=0{,}9\,u_{n}+120$. Soit la suite $(v_{n})$ définie pour tout $n\in\N$ par
$v_{n}=u_{n}-1\,200$.
\begin{questions}[label=\alph*)]
  \item Démontrer que la suite $(v_{n})$ est géométrique ; préciser sa raison et son
        premier terme.
  \item Exprimer $u_{n}$ en fonction de $n$.
  \item Déterminer la limite de la suite $(u_{n})$ et l'interpréter.
  \item À l'aide de la calculatrice, déterminer à partir de combien de mois le parc compte
        plus de $1\,100$ vélos.
\end{questions}
\begin{correction}
\cas{a)}
\explique{Soit $n\in\N$.}{}
\begin{calculs}
v_{n+1} &= u_{n+1}-1\,200 &\qquad& \rem{on remplace $n$ par $n+1$}\\
v_{n+1} &= 0{,}9\,u_{n}+120-1\,200 &\qquad& \rem{relation de récurrence}\\
v_{n+1} &= 0{,}9\,u_{n}-1\,080 &&\\
v_{n+1} &= 0{,}9\left(u_{n}-1\,200\right) &\qquad& \rem{$0{,}9\times 1\,200=1\,080$}\\
v_{n+1} &= 0{,}9\,v_{n} &&
\end{calculs}
\begin{calculs}
v_{0} &= 800-1\,200 = -400 &&
\end{calculs}
\explique{La suite $(v_{n})$ est géométrique de raison $\res{0{,}9}$ et de premier terme
$v_{0}=\res{-400}$.}{}
\cas{b)}
\begin{calculs}
v_{n} &= -400\times 0{,}9^{n} &\qquad& \rem{$v_{n}=v_{0}\times q^{n}$}\\
u_{n} &= v_{n}+1\,200 &\qquad& \rem{car $v_{n}=u_{n}-1\,200$}\\
u_{n} &= \res{1\,200-400\times 0{,}9^{n}} &&
\end{calculs}
\cas{c)}
\explique{$-1<0{,}9<1$, donc $\lim\limits_{n\to+\infty}0{,}9^{n}=0$ ; par produit puis par
somme, \res{$\lim\limits_{n\to+\infty}u_{n}=1\,200$}.}{}
\explique{À long terme, le parc se stabilise autour de $1\,200$ vélos.}{on revient au
contexte}
\cas{d)}
\begin{calculs}
u_{n} &> 1\,100 &&\\
1\,200-400\times 0{,}9^{n} &> 1\,100 &\qquad& \rem{on remplace $u_{n}$}\\
\iff -400\times 0{,}9^{n} &> -100 &\qquad& \rem{on retire $1\,200$}\\
\iff 0{,}9^{n} &< 0{,}25 &\qquad& \rem{on divise par $-400<0$ : l'ordre change}
\end{calculs}
\explique{À la calculatrice, $0{,}9^{13}\approx 0{,}254$ et $0{,}9^{14}\approx 0{,}229$ ; la
suite $\left(0{,}9^{n}\right)$ est décroissante, donc l'inégalité est vraie pour tout
$n\geq 14$.}{tableau de valeurs}
\explique{Le parc compte plus de $1\,100$ vélos \res{à partir de $14$ mois}.}{}
\end{correction}

\etape{À vous de jouer}
Un club d'astronomie compte $150$ membres en septembre. Chaque année, $20\,\%$ des membres
ne renouvellent pas leur adhésion, et $60$ nouveaux membres s'inscrivent. On note $u_{n}$ le
nombre de membres $n$ années après : $u_{0}=150$ et, pour tout $n\in\N$,
$u_{n+1}=0{,}8\,u_{n}+60$. Soit la suite $(w_{n})$ définie pour tout $n\in\N$ par
$w_{n}=u_{n}-300$.
\begin{questions}[label=\alph*)]
  \item Démontrer que la suite $(w_{n})$ est géométrique ; préciser sa raison et son
        premier terme.
  \item Exprimer $u_{n}$ en fonction de $n$.
  \item Déterminer la limite de la suite $(u_{n})$ et l'interpréter.
  \item À l'aide de la calculatrice, déterminer à partir de combien d'années le club
        compte plus de $280$ membres.
\end{questions}
\reponse{a) Raison $0{,}8$ et $w_{0}=-150$.\; b)~$u_{n}=300-150\times 0{,}8^{n}$.\;
c)~$300$.\; d)~$0{,}8^{n}<\frac{2}{15}$ : à partir de $10$ ans.}
\end{methodebox}

\afaire{exercices 34 à 37}

\partiecours{Théorème du point fixe}

Pour une suite $u_{n+1}=f(u_{n})$ qui converge, la limite est l'abscisse d'un point
d'intersection de la courbe de $f$ et de la droite d'équation $y=x$ (chapitre~1,
méthode~5). Le théorème qui le dit utilise la \textbf{continuité} de $f$, étudiée au
chapitre~6 ; pour l'instant, on en retient deux choses :
\begin{itemize}
  \item une fonction est continue sur un intervalle si sa courbe s'y trace « sans lever le
        crayon » ;
  \item une fonction dérivable sur un intervalle est continue sur cet intervalle : les
        polynômes, la fonction racine carrée sur $\intfo{0}{+\infty}$, les quotients sur
        un intervalle où le dénominateur ne s'annule pas sont continus.
\end{itemize}

\begin{proprietebox}[Théorème 4 --- Théorème du point fixe (admis)]
\begin{minipage}[c]{0.58\linewidth}\raggedright
Soit $f$ une fonction continue sur un intervalle $I$, et $(u_{n})$ une suite telle que,
pour tout $n\in\N$, $u_{n}\in I$ et $u_{n+1}=f(u_{n})$. Si $(u_{n})$ converge vers un
nombre réel $\ell$ de $I$, alors $f(\ell)=\ell$.

\smallskip
Ce théorème ne prouve pas la convergence : on la démontre d'abord (partie~VI). Il donne
seulement les valeurs possibles de $\ell$, solutions de l'équation $f(x)=x$.
\end{minipage}\hfill
\begin{minipage}[c]{0.38\linewidth}\centering
\begin{tikzpicture}[x=0.6cm,y=0.6cm]
  \repere{0}{0}{5}{5}
  \draw[black!70,line width=0.8pt] (0,0)--(5,5) node[above left,font=\footnotesize]{$y=x$};
  \draw[coulCourbe,line width=1.1pt] (0,2)--(5,4.5);
  \node[coulCourbe,font=\footnotesize] at (1,3) {$\Cf$};
  \draw[coulCourbeB,line width=0.7pt] (0,0)--(0,2)--(2,2)--(2,3)--(3,3)--(3,3.5)--(3.5,3.5);
  \draw[dashed,black!60] (4,4)--(4,0);
  \path[coulCourbeB] (0,0) node[croix]{} (2,0) node[croix]{} (3,0) node[croix]{}
    (3.5,0) node[croix]{};
  \path (4,4) node[croix]{};
  \node[font=\footnotesize,below,inner sep=2pt] at (4.1,0) {$\ell$};
\end{tikzpicture}
\end{minipage}
\end{proprietebox}

\begin{remarquebox}[Pourquoi ?]
Si $(u_{n})$ converge vers $\ell$, la suite $(u_{n+1})$, qui a les mêmes termes décalés
d'un rang, converge aussi vers $\ell$. Comme $f$ est continue, quand $u_{n}$ se
rapproche de $\ell$, $f(u_{n})$ se rapproche de $f(\ell)$. L'égalité $u_{n+1}=f(u_{n})$
devient, à la limite, $\ell=f(\ell)$.
\end{remarquebox}

\begin{methodebox}[Méthode 12 --- Déterminer la limite d'une suite $u_{n+1}=f(u_{n})$]
\etapeexemples
Soit $f$ la fonction définie sur $\intfo{0}{+\infty}$ par $f(x)=\dfrac{4x+3}{x+2}$ et la
suite $(w_{n})$ définie par $w_{0}=1$ et, pour tout $n\in\N$, $w_{n+1}=f(w_{n})$.
\begin{questions}[label=\alph*)]
  \item Étudier le sens de variation de $f$ sur $\intfo{0}{+\infty}$.
  \item Démontrer par récurrence que, pour tout $n\in\N$, on a
        $1\leq w_{n}\leq w_{n+1}\leq 3$.
  \item En déduire que la suite $(w_{n})$ converge.
  \item Déterminer la limite de $(w_{n})$.
\end{questions}
\begin{correction}
\cas{a)}
\explique{$f$ est un quotient dont le dénominateur $x+2$ ne s'annule pas sur
$\intfo{0}{+\infty}$ : elle y est dérivable. Soit $x\geq 0$.}{$u(x)=4x+3$ et $v(x)=x+2$}
\begin{calculs}
f'(x) &= \frac{4(x+2)-(4x+3)\times 1}{(x+2)^{2}} &\qquad& \rem{$\left(\frac{u}{v}\right)'=\frac{u'v-uv'}{v^{2}}$}\\
f'(x) &= \frac{4x+8-4x-3}{(x+2)^{2}} &\qquad& \rem{on développe}\\
f'(x) &= \frac{5}{(x+2)^{2}} &&
\end{calculs}
\explique{Pour tout $x\geq 0$, on a $f'(x)>0$ : \res{$f$ est strictement croissante sur
$\intfo{0}{+\infty}$}.}{un carré non nul est strictement positif}
\cas{b)}
\explique{Pour tout $n\in\N$, on note $P_{n}$ la propriété :
«~$1\leq w_{n}\leq w_{n+1}\leq 3$~».}{}
\explique{\textbf{Initialisation.} $w_{0}=1$ et $w_{1}=f(1)=\frac{7}{3}$ ;
$1\leq 1\leq\frac{7}{3}\leq 3$, donc $P_{0}$ est vraie.}{}
\explique{\textbf{Hérédité.} Soit $n\in\N$ tel que $P_{n}$ est vraie. Montrons que
$1\leq w_{n+1}\leq w_{n+2}\leq 3$.}{}
\begin{calculs}
1 &\leq w_{n}\leq w_{n+1}\leq 3 &\qquad& \rem{hypothèse de récurrence}\\
f(1) &\leq f(w_{n})\leq f(w_{n+1})\leq f(3) &\qquad& \rem{$f$ croissante : l'ordre est
conservé}\\
\frac{7}{3} &\leq w_{n+1}\leq w_{n+2}\leq 3 &\qquad& \rem{$f(3)=\frac{15}{5}$, soit $3$}
\end{calculs}
\explique{Comme $1\leq\frac{7}{3}$, on a $1\leq w_{n+1}\leq w_{n+2}\leq 3$ : $P_{n+1}$ est
vraie.}{on obtient mieux que le but}
\explique{\textbf{Conclusion.} $P_{0}$ est vraie et la propriété est héréditaire ;
d'après le principe de récurrence, \res{pour tout $n\in\N$, on a
$1\leq w_{n}\leq w_{n+1}\leq 3$}.}{}
\cas{c)}
\explique{Pour tout $n\in\N$, on a $w_{n}\leq w_{n+1}$ : $(w_{n})$ est croissante ; et
$w_{n}\leq 3$ : elle est majorée par $3$.}{}
\explique{D'après le théorème de convergence monotone, \res{$(w_{n})$ converge} vers un
nombre réel $\ell$. Comme $1\leq w_{n}\leq 3$ pour tout $n\in\N$, on a
$1\leq\ell\leq 3$.}{propriété~3}
\cas{d)}
\explique{$f$ est dérivable, donc continue, sur $\intfo{0}{+\infty}$ (question a) ; pour
tout $n\in\N$, $w_{n}\in\intfo{0}{+\infty}$, et $\ell\in\intfo{0}{+\infty}$. D'après le
théorème du point fixe, $f(\ell)=\ell$.}{on vérifie les hypothèses une à une}
\begin{calculs}
\frac{4\ell+3}{\ell+2} &= \ell &\qquad& \rem{équation du point fixe}\\
4\ell+3 &= \ell(\ell+2) &\qquad& \rem{on multiplie par $\ell+2\neq 0$}\\
\ell^{2}-2\ell-3 &= 0 &\qquad& \rem{une équation du second degré}
\end{calculs}
\begin{calculs}
\Delta &= (-2)^{2}-4\times 1\times(-3) &&\\
\Delta &= 16 &\qquad& \rem{$\Delta>0$ : deux solutions}\\
x_{1} &= \frac{2-4}{2} = -1 &&\\
x_{2} &= \frac{2+4}{2} = 3 &&
\end{calculs}
\explique{Or $1\leq\ell\leq 3$, donc $\ell\neq -1$ : \res{$(w_{n})$ converge vers
$3$}.}{l'encadrement élimine une solution}
\end{correction}

\etape{À vous de jouer}
Soit $g$ la fonction définie sur $\intfo{0}{+\infty}$ par $g(x)=\dfrac{6x+4}{x+3}$ et la
suite $(t_{n})$ définie par $t_{0}=0$ et, pour tout $n\in\N$, $t_{n+1}=g(t_{n})$.
\begin{questions}[label=\alph*)]
  \item Étudier le sens de variation de $g$ sur $\intfo{0}{+\infty}$.
  \item Démontrer par récurrence que, pour tout $n\in\N$, on a
        $0\leq t_{n}\leq t_{n+1}\leq 4$.
  \item En déduire que la suite $(t_{n})$ converge.
  \item Déterminer la limite de $(t_{n})$.
\end{questions}
\reponse{a) $g'(x)=\frac{14}{(x+3)^{2}}>0$ : $g$ croissante.\; b)~$t_{1}=\frac{4}{3}$ ;
$g(0)=\frac{4}{3}$ et $g(4)=4$.\; c)~Croissante et majorée par $4$.\; d)~$\ell^{2}-3\ell-4=0$,
solutions $-1$ et $4$ ; limite $4$.}
\end{methodebox}

\begin{remarquebox}[Attention]
Sans la convergence, l'équation $f(x)=x$ ne prouve rien. La suite $(z_{n})$ définie par
$z_{0}=2$ et, pour tout $n\in\N$, $z_{n+1}=2z_{n}-1$ vaut $2$ ; $3$ ; $5$ ; $9$… et tend vers $+\infty$, alors
que l'équation $x=2x-1$ a pour solution $1$.
\end{remarquebox}

\afaire{exercices 38 à 41, puis les exercices type bac 42 à 44}

\end{document}
